\documentclass[letterpaper]{article}
\usepackage{geometry}
\usepackage{xcolor}
\usepackage{amsmath}
\usepackage[some]{background}
\usepackage{array}
\usepackage{setspace}
\usepackage{tikz}
\usepackage[hidelinks]{hyperref}

\addtolength{\oddsidemargin}{-.875in}
\addtolength{\evensidemargin}{-.875in}
\addtolength{\textwidth}{1.75in}

\addtolength{\topmargin}{-.875in}
\addtolength{\textheight}{1.75in}

\definecolor{titlepagecolor}{cmyk}{1,.60,0,.40}
\definecolor{workheadingcolor}{cmyk}{0,.68,.73,.05}
\definecolor{sectionheadingcolor}{cmyk}{1,.60,0,.40}
% \definecolor{sectionheadingcolor}{cmyk}{0,.67,0.98,.01}

\newcolumntype{M}[1]{>{\centering\arraybackslash}m{#1}}
\newcolumntype{L}[1]{>{\arraybackslash}m{#1}}
\newcolumntype{R}[1]{>{\raggedleft\arraybackslash}m{#1}}

\renewcommand\fbox{\fcolorbox{titlepagecolor}{titlepagecolor}}

% Source: http://tex.stackexchange.com/questions/89588/positioning-relative-to-page-in-tikz
% Defining a new coordinate system for the page:
%
% --------------------------
% |(-1,1)    (0,1)    (1,1)|
% |                        |
% |(-1,0)    (0,0)    (1,0)|
% |                        |
% |(-1,-1)   (0,-1)  (1,-1)|
% --------------------------
\makeatletter
\def\parsecomma#1,#2\endparsecomma{\def\page@x{#1}\def\page@y{#2}}
\tikzdeclarecoordinatesystem{page}{
    \parsecomma#1\endparsecomma
    \pgfpointanchor{current page}{north east}
    % Save the upper right corner
    \pgf@xc=\pgf@x%
    \pgf@yc=\pgf@y%
    % save the lower left corner
    \pgfpointanchor{current page}{south west}
    \pgf@xb=\pgf@x%
    \pgf@yb=\pgf@y%
    % Transform to the correct placement
    \pgfmathparse{(\pgf@xc-\pgf@xb)/2.*\page@x+(\pgf@xc+\pgf@xb)/2.}
    \expandafter\pgf@x\expandafter=\pgfmathresult pt
    \pgfmathparse{(\pgf@yc-\pgf@yb)/2.*\page@y+(\pgf@yc+\pgf@yb)/2.}
    \expandafter\pgf@y\expandafter=\pgfmathresult pt
}
\makeatother

\backgroundsetup{
scale=1,
angle=0,
opacity=1,
contents={\begin{tikzpicture}[remember picture,overlay]
 \path [fill=titlepagecolor] (page cs:-1,0.79) rectangle (page cs:1,1); 
 \path [fill=titlepagecolor] (page cs:-0.95,0.61) rectangle (page cs:-0.91,0.70);
 \path [fill=titlepagecolor] (page cs:-0.95,-0.025) rectangle (page cs:-0.91,0.065);
% \path [fill=titlepagecolor] (page cs:-0.95,-.395) rectangle (page cs:-0.91,-.48);
 \path [fill=titlepagecolor] (page cs:0.04,0.555) rectangle (page cs:0.08,0.70);
 \path [fill=titlepagecolor] (page cs:0.04,0.1) rectangle (page cs:0.08,0.21);
 \draw [color=white, very thick] (page cs:0.5,0.8)--(page cs:0.5,1);
\end{tikzpicture}}
}

\begin{document}

\begin{titlepage}
\BgThispage
\newgeometry{left=1cm,right=6cm,bottom=2cm}
\vspace*{-3.5cm}
\noindent
\begin{minipage}{0.4\linewidth}
	\vspace{0.4cm}
	\textcolor{white}{\Huge\textbf{\textsf{Siddhartha}}}\\
	\begin{spacing}{0.5}
	\textcolor{white}{\normalsize \textsf{Machine Learning / FPGA Engineer}}
	\end{spacing}
\end{minipage}\hspace{30pt}
\begin{minipage}{0.3\linewidth}
	\vspace{0cm}
	\begin{center}
        \begin{tikzpicture}
            \clip (0,0)  circle (1.25cm) ;
            % adjust this coordinate to move image
            \node[anchor=center] at (0,-0.2) {\includegraphics[width=3cm]{me_2019}};
        \end{tikzpicture}
	\end{center}
\end{minipage}\hspace{125pt}
\begin{minipage}{0.3\linewidth}
	\vspace{-0.1cm}
	\begin{flushright}
        %!TEX root= ./cv.tex

\begin{tabular}{M{30mm}M{8mm}}
\multicolumn{1}{r}{\textcolor{white}{\small \textsf{sidmontu@gmail.com}}} & \multicolumn{1}{l}{\fbox{\includegraphics[width=12px]{email_logo}}}\\
    \multicolumn{1}{r}{\textcolor{white}{\small \textsf{\href{https://www.linkedin.com/in/siddhartha-63249064/}{siddhartha-63249064}}}} & \multicolumn{1}{l}{\fbox{\includegraphics[height=12px]{linkedin_logo}}}\\
% \multicolumn{1}{r}{\textcolor{white}{\small \textsf{Singapore}}} & \includegraphics[height=10px]{map_logo} \\
\multicolumn{1}{r}{\textcolor{white}{\small \textsf{\href{https://www.github.com/sidmontu}{github.com/sidmontu}}}} & \multicolumn{1}{l}{\fbox{\includegraphics[height=12px]{github}}} \\
% \multicolumn{1}{r}{\textcolor{white}{\small \textsf{\href{https://www.twitter.com/sidmontu}{twitter.com/sidmontu}}}} & \multicolumn{1}{l}{\fbox{\includegraphics[height=8px]{twitter_logo}}}\\
\multicolumn{1}{r}{\textcolor{white}{\small \textsf{\href{https://sidmontu.github.io}{sidmontu.github.io}}}} & \multicolumn{1}{l}{\fbox{\includegraphics[height=12px]{globe_logo}}}
\end{tabular}

	\end{flushright}
\end{minipage}

\vspace*{0.6cm}

\noindent
    % \fbox{
        \begin{minipage}{0.5\paperwidth}
        %!TEX root= ./cv.tex

{\Large \textbf{\textcolor{sectionheadingcolor}{WORK EXPERIENCE}}}\\
\\
\begin{tabular}{L{0.1em} L{18em} L{5.9em}}

    \multicolumn{2}{l}{\raisebox{-.15\height}{\includegraphics[width=9px]{calendar_logo}} \textcolor{workheadingcolor}{\scriptsize\textbf{\textsf{Feb 2020 - Present}}}} & \multicolumn{1}{R{5.9em}}{\raisebox{-.25\height}{\includegraphics[width=9px]{location_icon}} \textcolor{workheadingcolor}{\scriptsize\textbf{\textsf{Singapore}}}}\\

    \multicolumn{2}{l}{\rule{0pt}{2.5ex}{\small \textsf{\bf Founder, CTO, Director}}} & \\
    \multicolumn{2}{l}{{\small\textsf{\color{titlepagecolor}\bf inPact.ai}}} & \\

    % \textcolor{workheadingcolor}{\scriptsize\textbf{\textsf{Accomplishments}}} & \\ 

    \raisebox{1.25\height}{\rule{0pt}{3ex}\tikz \fill [titlepagecolor] (-0.07,-0.2) rectangle (0.07,-0.05);} & \multicolumn{2}{L{23.9em}}{{\footnotesize \textsf{Led product development across all aspects of machine learning, frontend, backend, and deployment.}}}\\
    \raisebox{5\height}{\tikz \fill [titlepagecolor] (-0.07,-0.2) rectangle (0.07,-0.05);} & \multicolumn{2}{L{23.9em}}{{\footnotesize \textsf{Hired and managed contractors and an intern to meet product development targets regularly. Enforced good development practices such as version control, linters, code reviews, etc.}}}\\
    % \raisebox{2\height}{\tikz \fill [titlepagecolor] (-0.07,-0.2) rectangle (0.07,-0.05);} & \multicolumn{2}{L{23.9em}}{{\footnotesize \textsf{Released two major product iterations: the MVP {\it Alpha} release, and an improved {\it Beta} version based on customer feedback.}}}\\
    \raisebox{5\height}{\tikz \fill [titlepagecolor] (-0.07,-0.2) rectangle (0.07,-0.05);} & \multicolumn{2}{L{23.9em}}{{\footnotesize \textsf{Gained hands-on development experience with MLOps tools: Label studio for data annotation, MLFlow for model registry, Hub for dataset/feature store, and Cortex.dev for model deployment.}}}\\
    \raisebox{5\height}{\tikz \fill [titlepagecolor] (-0.07,-0.2) rectangle (0.07,-0.05);} & \multicolumn{2}{L{23.9em}}{{\footnotesize \textsf{Hands-on experience with ReactJS and serverless frameworks. Setup inPact's cloud infrastructure on AWS -- used a myriad of AWS cloud services such as Cognito, DynamoDB, S3, SNS, Lambda, etc.}}}\\
    \raisebox{2\height}{\tikz \fill [titlepagecolor] (-0.07,-0.2) rectangle (0.07,-0.05);} & \multicolumn{2}{L{23.9em}}{{\footnotesize \textsf{Handled business functions such as customer acquisition, corporate partnerships, investor relations, fundraising, marketing, and more.}}}\\
    \raisebox{5\height}{\tikz \fill [titlepagecolor] (-0.07,-0.2) rectangle (0.07,-0.05);} & \multicolumn{2}{L{23.9em}}{{\footnotesize \textsf{Raised S\$75K pre-seed from Entrepreneur First, an international VC firm funded by Reid Hoffman (founder of LinkedIn), founders of DeepMind and PayPal, and some of the top investors in the world.}}}\\
    
    % & &\\

    \multicolumn{2}{l}{\rule{0pt}{5ex}\raisebox{-.15\height}{\includegraphics[width=9px]{calendar_logo}} \textcolor{workheadingcolor}{\scriptsize\textbf{\textsf{Oct 2017 - Dec 2019}}}} & \multicolumn{1}{r}{\raisebox{-.25\height}{\includegraphics[width=9px]{location_icon}} \textcolor{workheadingcolor}{\scriptsize\textbf{\textsf{Sydney, Australia}}}}\\

    \multicolumn{2}{l}{\rule{0pt}{2.5ex}{\small \textsf{\bf Postdoctoral Research Associate}}} & \\
    \multicolumn{2}{l}{{\small\textsf{\color{titlepagecolor}\bf University of Sydney}}} & \\

    % \textcolor{workheadingcolor}{\scriptsize\textbf{\textsf{Accomplishments}}} & \\ 

    \raisebox{2.25\height}{\rule{0pt}{3ex}\tikz \fill [titlepagecolor] (-0.07,-0.2) rectangle (0.07,-0.05);} & \multicolumn{2}{L{23.9em}}{{\footnotesize \textsf{Conducted research on: next-generation FPGA (overlay) architectures, low-precision deep neural networks, on-chip machine learning, and RF communication systems (Ettus RFNoC framework).}}}\\
    \raisebox{2.25\height}{\tikz \fill [titlepagecolor] (-0.07,-0.2) rectangle (0.07,-0.05);} & \multicolumn{2}{L{23.9em}}{{\footnotesize \textsf{Core team member on a project (High-Speed Machine Learning for RF Communication) commissioned by the Australian Defense Force.}}}\\
    \raisebox{0\height}{\tikz \fill [titlepagecolor] (-0.07,-0.2) rectangle (0.07,-0.05);} & \multicolumn{2}{L{23.9em}}{{\footnotesize \textsf{(Co-)Authored 5 research papers/posters during the stint at the lab.}}}\\
    \raisebox{1.75\height}{\tikz \fill [titlepagecolor] (-0.07,-0.2) rectangle (0.07,-0.05);} & \multicolumn{2}{L{23.9em}}{{\footnotesize \textsf{Supervised final year undergraduate projects, assisted with teaching/invigilation, and undertook sysadmin duties over lab resources.}}}\\


%     \textcolor{workheadingcolor}{\scriptsize\textbf{\textsf{Oct 2017 - Dec 2019}}} & \textcolor{workheadingcolor}{\scriptsize\textbf{\textsf{Sydney, Australia}}}\\
% 
%     {\small \textbf{\textsf{Postdoctoral Research Associate}}} & \\
%     {\small\textsf{University of Sydney}} & \\
% 
%     % \textcolor{workheadingcolor}{\scriptsize\textbf{\textsf{Accomplishments}}} & \\ 
% 
%     {\footnotesize \textsf{My research projects have taught me a variety of hardware and low-level embedded design techniques, and have enabled me to build a strong understanding of computer architecture.}} &\\
% 
%     &\\
% 
%     \textcolor{workheadingcolor}{\scriptsize\textbf{\textsf{Jun 2012 - Oct 2012}}} & \textcolor{workheadingcolor}{\scriptsize\textbf{\textsf{London, UK}}}\\
% 
%     {\small \textbf{\textsf{Summer Research Internship}}} & \\
%     {\small\textsf{Imperial College London}} & \\
% 
%     % \textcolor{workheadingcolor}{\scriptsize\textbf{\textsf{Accomplishments}}} & \\ 
% 
%     {\footnotesize \textsf{Summer research project on abcd. Learnt wxyz, and achieved jklm.}} &\\
% 
%     &\\

\end{tabular}

        \vspace{0.2in}
        %!TEX root= ./cv.tex

{\Large \textbf{\textcolor{sectionheadingcolor}{SKILLS \& COMPETENCIES}}}\\
\\
\begin{tabular}{L{0.1em} L{18em} L{5.9em}}

    \multicolumn{3}{l}{PyTorch / Tensorflow} \\
    \multicolumn{3}{l}{C / C++} \\
    \multicolumn{3}{l}{Verilog / VHDL} \\
    \multicolumn{3}{l}{Typesetting in \LaTeX} \\
    \multicolumn{3}{l}{PyTorch / Tensorflow} \\
    \multicolumn{3}{l}{PyTorch / Tensorflow} \\
    \multicolumn{3}{l}{PyTorch / Tensorflow} \\
    \multicolumn{3}{l}{PyTorch / Tensorflow} \\
    \multicolumn{3}{l}{PyTorch / Tensorflow} \\
    \multicolumn{3}{l}{PyTorch / Tensorflow} \\
    \multicolumn{3}{l}{PyTorch / Tensorflow} \\
    \multicolumn{3}{l}{PyTorch / Tensorflow} \\
    \multicolumn{3}{l}{PyTorch / Tensorflow} \\
    \multicolumn{3}{l}{PyTorch / Tensorflow} \\
    \multicolumn{3}{l}{PyTorch / Tensorflow} \\


\end{tabular}


        \end{minipage}\hspace{-5pt}
    % }
%
\begin{minipage}{0.5\paperwidth}
	\vspace{0cm}
    %!TEX root= ./cv.tex

{\Large \textbf{\textcolor{sectionheadingcolor}{EDUCATION}}}\\
\\
\begin{tabular}{L{0.1em} L{18em} L{5.9em}}

    % PhD NTU
    \multicolumn{2}{L{18.1em}}{\raisebox{-.15\height}{\includegraphics[width=9px]{calendar_logo}} \textcolor{workheadingcolor}{\scriptsize\textbf{\textsf{Jan 2013 - Jan 2019}}}} & \multicolumn{1}{R{5.9em}}{\raisebox{-.25\height}{\includegraphics[width=9px]{location_icon}} \textcolor{workheadingcolor}{\scriptsize\textbf{\textsf{Singapore}}}}\\

    \multicolumn{3}{l}{\rule{0pt}{1.5ex}{\small \textsf{\bf PhD, Computer Science \& Engineering}}} \\
    \multicolumn{3}{l}{{\small\textsf{\color{titlepagecolor}\bf Nanyang Technological University (NTU)}}} \\

    \multicolumn{3}{l}{{\footnotesize \textsf{Dissertation Title: \href{https://sidmontu.github.io/thesis.pdf}{\textcolor{titlepagecolor}{Dataflow Optimized Overlays for FPGAs}}}}}\\
    \multicolumn{3}{l}{{\footnotesize\textsf{Supervisor: \href{https://nachiket.github.io}{\textcolor{titlepagecolor}{Dr. Nachiket Kapre}}}}}\\
    \raisebox{1.25\height}{\rule{0pt}{3ex}\tikz \fill [titlepagecolor] (-0.07,-0.2) rectangle (0.07,-0.05);} & \multicolumn{2}{L{23.9em}}{{\footnotesize \textsf{Introduces DaCO, a Dataflow Coprocess Overlay optimized for the Arria 10 FPGAs. Research challenges addressed include:}}}\\
    \raisebox{1.25\height}{\tikz \fill [titlepagecolor] (-0.07,-0.2) rectangle (0.07,-0.05);} & \multicolumn{2}{L{23.9em}}{{\footnotesize \textsf{Custom dataflow scheduling circuit that enables large-scale out-of-order instruction execution at runtime,}}}\\
    \raisebox{1.25\height}{\tikz \fill [titlepagecolor] (-0.07,-0.2) rectangle (0.07,-0.05);} & \multicolumn{2}{L{23.9em}}{{\footnotesize \textsf{Priority-aware NoC packet routing for improved criticality-aware dataflow communication, and}}}\\
    \raisebox{1.25\height}{\tikz \fill [titlepagecolor] (-0.07,-0.2) rectangle (0.07,-0.05);} & \multicolumn{2}{L{23.9em}}{{\footnotesize \textsf{Compiler support that optimizes the dataflow graph for improved runtime execution.}}}\\
    \raisebox{1.25\height}{\tikz \fill [titlepagecolor] (-0.07,-0.2) rectangle (0.07,-0.05);} & \multicolumn{2}{L{23.9em}}{{\footnotesize \textsf{(Co-)Authored a total of 15 research papers/posters/journals at top-tier IEEE conferences during the candidacy.}}}\\

    % BEng, Imperial College London
    \multicolumn{2}{L{18.1em}}{\rule{0pt}{3ex}\raisebox{-.15\height}{\includegraphics[width=9px]{calendar_logo}} \textcolor{workheadingcolor}{\scriptsize\textbf{\textsf{Oct 2009 - Jun 2012}}}} & \multicolumn{1}{R{5.9em}}{\raisebox{-.25\height}{\includegraphics[width=9px]{location_icon}} \textcolor{workheadingcolor}{\scriptsize\textbf{\textsf{London, UK}}}}\\

    \multicolumn{3}{l}{\rule{0pt}{1.5ex}{\small \textsf{\bf Bachelors of Engineering (BEng)}}} \\
    \multicolumn{3}{l}{{\small \textsf{Electrical \& Electronics Engineering (EEE)}}} \\
    \multicolumn{3}{l}{{\small\textsf{\color{titlepagecolor}\bf Imperial College London}}} \\

\end{tabular}


    \vspace{0.2in}
    %!TEX root= ./cv.tex

{\Large \textbf{\textcolor{sectionheadingcolor}{NOTABLE PUBLICATIONS}}}\\
\\
\begin{tabular}{L{0.1em} L{18em} L{5.9em}}

    \raisebox{4.5\height}{\tikz \fill [titlepagecolor] (-0.07,-0.2) rectangle (0.07,-0.05);} & \multicolumn{2}{L{23.9em}}{{\footnotesize \textsf{\bf DaCO: A High-Performance Token Dataflow Coprocessor Overlay for FPGAs}, {\it Siddhartha, Nachiket Kapre}, International Conference on Field-Programmable Technology 2018}}\\
    \raisebox{4.5\height}{\tikz \fill [titlepagecolor] (-0.07,-0.2) rectangle (0.07,-0.05);} & \multicolumn{2}{L{23.9em}}{{\footnotesize \textsf{\bf Hoplite-Q: Priority-Aware Routing in FPGA Overlay NoCs}, {\it Siddhartha, Nachiket Kapre}, IEEE International Symposium on Field-Programmable Custom Computing Machines 2018}}\\
    \raisebox{7.5\height}{\tikz \fill [titlepagecolor] (-0.07,-0.2) rectangle (0.07,-0.05);} & \multicolumn{2}{L{23.9em}}{{\footnotesize \textsf{\bf LUXOR: An FPGA Logic Cell Architecture for Eficient Compressor Tree Implementations}, {\it Seyedramin Rasoulinezhad, Siddhartha, Hao Zhou, Lingli Wang, David Boland, Philip Leong}, ACM International Symposium on FPGAs 2020}}\\
    \raisebox{10\height}{\tikz \fill [titlepagecolor] (-0.07,-0.2) rectangle (0.07,-0.05);} & \multicolumn{2}{L{23.9em}}{{\footnotesize \textsf{\bf Long Short-Term Memory for Radio Frequency Spectral Prediction and its Real-Time FPGA Implementation}, {\it Siddhartha, Yee Hui Lee, Duncan Moss, Julian Faraone, Perry Blackmore, Daniel Salmond, David Boland, and Philip Leong}, IEEE Military Communications Conference 2018}}\\
    \raisebox{7.5\height}{\tikz \fill [titlepagecolor] (-0.07,-0.2) rectangle (0.07,-0.05);} & \multicolumn{2}{L{23.9em}}{{\footnotesize \textsf{\bf Real-Time Automatic Modulation Classification using RFSoC}, {\it Stephen Tridgell, David Boland, Philip Leong, Ryan Kastner, Alireza Khodamoradi, Siddhartha}, 27th Reconfigurable Architectures Workshop 2020}}\\

\end{tabular}


    \vspace{0.2in}
    %!TEX root= ./cv.tex

{\Large \textbf{\textcolor{sectionheadingcolor}{INTERESTS}}}\\
\\
\begin{tabular}{L{0.1em} L{18em} L{5.9em}}

    \multicolumn{3}{l}{Applied AI} \\
    \multicolumn{3}{l}{Computer Vision} \\
    \multicolumn{3}{l}{Natural Language Processing} \\
    \multicolumn{3}{l}{Computer Architecture} \\
    \multicolumn{3}{l}{FPGA Technologies} \\
    \multicolumn{3}{l}{RF Communication Systems} \\
    % \multicolumn{3}{l}{High-Performance Computing} \\
    % \multicolumn{3}{l}{Entrepreneurship} \\

\end{tabular}


\end{minipage}
\end{titlepage}
\restoregeometry
\end{document}
